\par\medskip\pdfbookmark[1]{Zu III und IV.}{nexus-note-III-IV}\centerline{Zu III und IV.}\smallskip
Ist $C$ der Complex aller positiven, nicht eigentlich-äquivalenten formae proprie primitivae von negativem Determinant $-D$, und legt man den Variabeln dieser Formen je zwei Werthe bei, welche relative Primzahlen zu einander sind, so ist die Anzahl aller Darstellungen einer positiven ganzen Zahl $m$ gleich $\epsilon\psi(m)$, wo $\epsilon$ die Anzahl der Auflösungen der Gleichung $tt+Duu=1$, und $\psi(m)$ die Anzahl derjenigen Wurzeln $n$ der Congruenz $nn+D\equiv0\pmod m$ bedeutet, für welche die drei Zahlen $m,\ 2n$ und $\frac{nn+D}m$ ohne gemeinschaftlichen Divisor sind (Disqq. Arithm. art. 180). Der Factor $\epsilon$ ist $=4$ für $D=1$, in allen andern Fällen $=2$. Ist ferner $m=p^\pi p'^{\pi'}p''^{\pi''}\ldots$, wo $p,p',p''\ldots$ von einander verschiedene Primzahlen bedeuten, so ist $\psi(m)=\psi(p^\pi)\psi(p'^{\pi'})\psi(p''^{\pi''})\ldots$; bedeutet $\mathfrak A(m)$ die Anzahl der Wurzeln $n$ der Congruenz $nn+D\equiv0\pmod m$, und bedient man sich des von \textsc{Legendre} eingeführten, von \textsc{Jacobi} verallgemeinerten Zeichens, so ist $\psi(p^\pi)=\mathfrak A(p^\pi)=1+\left(\frac{-D}p\right)$, wenn $p$ nicht in $2D$ aufgeht, sonst aber $=\mathfrak A(p^\pi-\frac1p\mathfrak A(p^{\pi+1})$; die Anzahl $\mathfrak A(p^\pi)$ lässt sich immer leicht bestimmen (Disqq. Arithm. art. 104), für die Folge reicht aber die Bemerkung aus, dass $\mathfrak A(p^\pi)$ immer von $\pi$ unabhängig wird, sobald $\pi$ eine gewisse Grösse überschreitet.

Legt man den Variabeln der in dem Complex $C$ enthaltenen Formen alle ganzzahligen Werthe ohne Ausnahme bei (Disqq. Arithm. art. 181), so wird die Anzahl $(m)$ aller Darstellungen der Zahl $m$ gleich $\epsilon f(m)$, wo $f(m)=\Sigma\psi\left(\frac m{\mu\mu}\right)$ ist, und das Summenzeichen sich auf alle quadratischen Divisoren $\mu\mu$ der Zahl $m$ bezieht. Hieraus folgt unmittelbar
\[
f(m)=f(p^\pi p'^{\pi'}p''^{\pi''}\ldots)=f(p^\pi)f(p'^{\pi'})f(p''^{\pi''})\ldots
\]
und
\[
f(p^\pi)=\psi(p^\pi)+\psi(p^{\pi-2})+\psi(p^{\pi-4})+\ldots
\]
welche Reihe so lange fortzusetzen ist, als die Exponenten $\pi,\ \pi-2,\ \pi-4\ldots$ nicht negativ werden. Wenn $p$ nicht in $2D$ aufgeht, so folgt hieraus
\[
f(p^\pi)=1+\left(\frac{-D}p\right)+\left(\frac{-D}{pp}\right)+\ldots+\left(\frac{-D}{p^\pi}\right)
\]
und allgemein, wenn $m$ relative Primzahl zu $2D$ ist,
\[
f(m)=\Sigma\left(\frac{-D}n\right)
\]
wo das Summenzeichen sich auf alle Divisoren $n$ der Zahl $m$ bezieht.

\SourceLeaf{304}
Aus diesen Bemerkungen ergiebt sich unmittelbar die Richtigkeit der im Text (III, 1, 2, 3) aufgestellten Sätze über die Anzahl $(m)$, wenn man für den ersten derselben noch die Bedingung hinzufügt, dass $D$ nicht durch $pp$ theilbar sein darf (die Bestimmung der Classenanzahl ist schon in den Disqq. Arithm. art. 256 auf den Fall zurückgeführt, in welchem $D$ durch kein Quadrat theilbar ist). Zugleich findet man, auch ohne Rücksicht auf diese Beschränkung, dass die unendliche Reihe
\[
1+\frac{f(p)}p+\frac{f(pp)}{pp}+\frac{f(p^3)}{p^3}+\ldots
\]
den Werth
\[
\frac1{1-\frac1p}\quad\text{oder}\quad\frac1{1-.\frac1p}\pdot\frac1{1-\left(\frac{-D}p\right)\frac1p}
\]
hat, je nachdem $2D$ durch die Primzahl $p$ theilbar oder nicht theilbar ist.
\BookRule{.25\textwidth}
